\documentclass[text06.tex]{subfiles}
\begin{document}
\newpage
\section{まとめ}
Open JTalkとJuliusの使いかたをまとめました。

Open JTalkの使いかたは、HSPスクリプトエディタで\\
\code{jtload “<読み上げたい文字列>”, <バッファ番号>}\\
で音声をロードした後、\\
\code{mmplay <バッファ番号>}\\
によって音声を再生します。\\

Juliusは、この教材で用意しているコマンドを使います。ターミナルでは\\
\code{julius.sh <.dicファイル>}\\
のように使います。人が編集しやすい \code{.yomi} ファイル（識別子（コンピュータが使うラベル）と読み方のペア）を作って、ターミナルに\\
\mbox{\code{convert\_yomi.sh}} (今作った.yomiファイル) > (出力するファイルの名前(.dicファイル))
とすると、Juliusが利用する \code{.dic} ファイルに変換できます。

HSPから使うときは、\code{julius.as} がJuliusとHSPの間で認識結果を受け渡してくれます。まず、特に大事なのは次の4つです。\\
\code{init\_julius <dicファイル>}\\
で辞書ファイルを指定しつつ、音声認識を使う準備をします。\\
\code{get\_word\_list words, cm}\\
で認識結果を受け取ります。\code{words(cnt)} に認識した言葉、\code{cm(cnt)} に自信の目安が入ります。\code{cm} は1に近いほど自信が高い結果ですが、正解する\ruby{確率}{かく|りつ}そのものではありません。\\
プログラムを読むときは \code{is\_recieved(sockidx)} も使います。Juliusから認識結果が届いたかを調べ、届いていれば0が返ってきます。

\subsection{\ruby{発展}{はっ|てん}:文法を設定する}
複数の言葉を組み合わせて認識したいときには、「\ruby{文法}{ぶん|ぽう}」という仕組みを使う方法もあります。

例えば、1から1000までの数字を一つずつ単語として登録するのは大変です。「ひゃく」「にじゅう」「さん」のような言葉を用意して、それらをどのように組み合わせられるかを「文法」で決める方法があります。「ひゃく」＋「にじゅう」＋「さん」→「ひゃくにじゅうさん」のように組み合わせることができます。7には「なな」や「しち」のように複数の読み方があることもあります。興味があれば、先生に聞いてみましょう。
\end{document}
